\chapter{Introduction}

Textile manufacturing is among the oldest and largest industrial sectors, yet its quality-control practices remain heavily dependent on human visual inspection. This chapter introduces the problem of automated fabric-defect detection, surveys the relevant literature, and establishes the motivation, problem statement, and objectives of the UltraFabric-Vision project. It closes with a brief description of the methodology, the assumptions and constraints under which the work is valid, and the organization of the remainder of the report.

\section[Introduction]{\textbf{Introduction}}

Woven and knitted fabric is produced on high-speed looms that generate cloth at roughly 30 metres per minute. At that rate, weaving faults, stains, holes, and structural irregularities can accumulate across long rolls before an operator notices them, turning a small fault into a large financial loss. Manual inspection is fatiguing and inconsistent: reported detection rates lie between 60\% and 75\%, and inspectors cannot sustain attention at line speed~\cite{kumar2019survey}. Automated visual inspection therefore offers a direct economic benefit, provided it can operate in real time and adapt to the enormous diversity of fabrics found on a factory floor.

Classical machine-vision approaches to fabric inspection relied on hand-crafted texture descriptors and supervised classifiers. These require large labelled datasets of defective samples, which are expensive to collect and quickly become obsolete as new defect types appear. UltraFabric-Vision instead adopts an \emph{anomaly-detection} formulation: the system learns only what defect-free fabric looks like and flags any statistically significant deviation. This removes the dependence on labelled defects and generalizes naturally to unseen fault types.

The project fuses three complementary deep-learning detectors---PatchCore~\cite{roth2022patchcore}, a self-supervised DINO Vision Transformer~\cite{caron2021dino}, and a Vision Transformer Autoencoder~\cite{dosovitskiy2021vit}---into a single calibrated ensemble, and packages it as a deployable, real-time inspection service.

\section[Literature Review]{\textbf{Literature Review}}

Automated fabric inspection has evolved from fixed texture descriptors to learned representations. Kumar~\cite{kumar2019survey} surveys computer-vision fabric-defect methods and notes that hand-crafted features (Gabor filters, co-occurrence statistics, local binary patterns) saturate at 70--85\% accuracy because they cannot adapt to new textures. Li et al.~\cite{li2020yolo} adapted YOLOv3 with fabric-specific anchor boxes and reported 89.7\% mean average precision at 45~FPS, while Jing et al.~\cite{jing2022mobileunet} proposed Mobile-Unet, a lightweight encoder--decoder achieving 96.3\% pixel-level segmentation on denim defects. These supervised methods perform well but depend on curated defect datasets.

Industrial anomaly detection reframed the problem around normal-only training. The MVTec~AD benchmark~\cite{bergmann2019mvtec} standardized evaluation across fifteen categories. PatchCore~\cite{roth2022patchcore} builds a coreset memory bank of normal patch features and scores anomalies by nearest-neighbour distance, achieving state-of-the-art localization; PaDiM~\cite{defard2021padim} models patch features with multivariate Gaussians. In parallel, self-supervised Vision Transformers such as DINO~\cite{caron2021dino} were shown to learn semantic structure in their attention maps without any labels, a property directly useful for separating coherent texture from anomalies. The Vision Transformer~\cite{dosovitskiy2021vit}, built on the self-attention mechanism of Vaswani et al.~\cite{vaswani2017attention}, provides the backbone for both feature extraction and reconstruction-based scoring. Public fabric datasets such as AITEX~\cite{silva2020aitex} support benchmarking, though real production data remains scarce.

The consensus from this literature is that (i) normal-only anomaly detection is the practical choice for evolving defect types, and (ii) different detection mechanisms capture different defect classes. What the literature under-emphasizes---and what this project addresses---is the engineering discipline required to make such an ensemble correct and fast in deployment: consistent preprocessing, score calibration across heterogeneous detectors, and accelerator-native scoring.

\section[Motivation]{\textbf{Motivation}}

Small and medium textile enterprises rarely have access to proprietary commercial inspection systems, which are expensive and closed. An open, self-supervised framework that can be calibrated to a factory's own fabric with a short pass over defect-free samples---and that runs at line speed on a single consumer GPU---would substantially lower the barrier to automated quality control. The motivation for this work is both to build such a system and to establish the engineering principles (preprocessing parity, score calibration, GPU-native scoring) that make it reliable enough to deploy.

\section[Problem statement]{\textbf{Problem statement}}

To design and industrialize a real-time, unsupervised fabric-defect detection system that fuses multiple complementary anomaly detectors into a single calibrated decision, achieves accurate discrimination between normal and defective cloth, and delivers low-latency inference suitable for integration with high-speed textile production lines.

\section[Objectives]{\textbf{Objectives}}
The objectives of the project are
\begin{enumerate}
\item To design an ensemble anomaly-detection pipeline that fuses PatchCore, DINO, and a Vision Transformer Autoencoder into a single, calibrated decision score.
\item To minimize inference latency through accelerator-native nearest-neighbour search and mixed-precision execution, enabling real-time operation.
\item To industrialize the system as a deployable, configurable service with a documented integration interface and validated end-to-end behaviour.
\end{enumerate}

\section[Brief Methodology of the project]{\textbf{Brief Methodology of the project}}

The methodology proceeds in four stages. First, defect-free fabric images are used to fit the two memory banks (PatchCore, DINO) and to train the reconstruction autoencoder, under a single preprocessing routine. Second, a calibration pass over the same normal data records each detector's score statistics, which are used to standardize the three heterogeneous scores to a common scale before fusion. Third, the fused decision threshold is fixed by maximizing Youden's~J statistic on a held-out validation set. Finally, the calibrated pipeline is wrapped in a REST/WebSocket service, an embeddable SDK, and a containerized deployment, and validated by an end-to-end test suite. The detailed methodology is developed in Chapters~3 and~4.

\section[Assumptions made / Constraints of the project]{\textbf{Assumptions made / Constraints of the project}}

The work is valid under the following assumptions and constraints. (i) The training set consists of \emph{defect-free} fabric that is representative of the production fabric; the calibrated threshold is only valid for the fabric family it was calibrated on. (ii) Imaging conditions---camera geometry, focus, and illumination---are held reasonably constant between calibration and inference; large changes require recalibration. (iii) Real-time latency figures assume a CUDA-capable GPU; on CPU the system is suitable only for offline batch inspection. (iv) The quantitative results reported here were obtained on a controlled synthetic fabric dataset; validation on real production fabric is required before field deployment.

\section[Organization of the report]{\textbf{Organization of the report}}

This report is organized as follows.
\begin{itemize}
\item Chapter 2 develops the theory and fundamentals: anomaly detection, the Vision Transformer, and the three detector pathways, together with the tools used.
\item Chapter 3 presents the system design---the parallel-pathway architecture, the calibration and fusion equations, and the deployment design.
\item Chapter 4 details the implementation, including the preprocessing-parity fix, GPU-accelerated scoring, calibration, and the service stack.
\item Chapter 5 reports and discusses the experimental results.
\item Chapter 6 concludes the report and outlines the future scope.
\end{itemize}
